\section{Results}

We demonstrate the five instrument properties in the order a user meets them, on
the fitted eight-axis map of $N=9{,}013$ patients. A single real patient
(\S\ref{sec:res-project}) makes the whole chain concrete; the remaining
subsections establish each property at the level of the full sample. Throughout,
where a specific axis or patient appears it illustrates a measurement property, not
a clinical claim.

\begin{figure}[t]
  \centering
  \includegraphics[width=\linewidth]{fig1_projection.png}
  \caption{\textbf{Projection --- the instrument returns a coordinate, its
  uncertainty, a summary, a prognosis, and the next measurement, for one patient.}
  A real patient from the bipolar cohort projected through the measurement model
  (Eq.~\ref{eq:m-scores}). \textbf{(A)} Posterior coordinate on all eight axes with
  $95\%$ credible intervals; axes with no observed indicator return the prior and
  are labelled as such, alongside the per-axis count of observed items.
  \textbf{(B)} The patient's archetype-simplex membership, concentrating $86\%$ of
  weight on one corner. \textbf{(C)} The interval-valued prognosis read from the
  coordinate's neighbourhood. \textbf{(D)} Value of information: the single next
  measurement that most reduces this patient's coordinate uncertainty, lifting the
  prior-dominated cognition axis to reliability $0.65$. The panels are one object
  read four ways.}
  \label{fig:projection}
\end{figure}

\subsection{Projection: a coordinate with its uncertainty, per patient}
\label{sec:res-project}

Figure~\ref{fig:projection} follows one patient from the bipolar cohort through the
instrument. Projecting their observed cells through \eqref{eq:m-scores} returns a
posterior coordinate on all eight axes together with its uncertainty
(Fig.~\ref{fig:projection}A). The two are inseparable and they are heterogeneous
across axes for the same patient: the immunometabolic coordinate is pinned down to
a posterior standard deviation of $0.22$ by six observed indicators, whereas
cognition and suicidality --- on which this patient has no observed item --- return
the prior, a standard deviation near $0.99$. The map does not paper over the gap
with a point estimate; it reports a wide interval and labels the axis
prior-dominated. This is the instrument's first and most basic honesty: a
coordinate is always accompanied by how much data produced it.

The same posterior supports three immediate reads. The patient's archetype mixture
(Fig.~\ref{fig:projection}B) places $86\%$ of their weight on a single corner of
the simplex, so the eight-number coordinate compresses to an interpretable
one-line summary without discarding the minority mixture. The interval-valued
prognosis (Fig.~\ref{fig:projection}C) is read off the coordinate's neighbourhood,
and --- because the coordinate carries $S_i$ --- it too is an interval, not a point.
And the value-of-information panel (Fig.~\ref{fig:projection}D) names the single
most informative next measurement for this patient: one verbal-memory test would
lift the cognition axis from prior-dominated to a reliability of $0.65$. Projection,
summary, prognosis and next-measurement all read off one object.

\begin{figure}[t]
  \centering
  \includegraphics[width=\linewidth]{fig2_information.png}
  \caption{\textbf{Information, not loading, localizes a patient.} For every
  modelled indicator, its loading on its home axis versus the Fisher information it
  contributes about that axis under its own likelihood family
  (Eq.~\ref{eq:m-fisher}). Continuous items convert loading into information
  efficiently; binary flags carry information $\lambda^2 p(1-p)$ that collapses at
  low endorsement, so a strongly-loading rare symptom localizes almost no one. The
  right panel translates this into a target-driven item count: reaching a posterior
  standard deviation of $0.5$ needs $\sim\!4$ items at loading $0.7$ but $\sim\!31$
  at loading $0.3$.}
  \label{fig:information}
\end{figure}

\subsection{Information: what localizes a patient is not the loading}
\label{sec:res-info}

Reading the map as an instrument reorders the item bank. Figure~\ref{fig:information}
contrasts each indicator's loading on its home axis with the Fisher information it
actually contributes under its own likelihood family \eqref{eq:m-fisher}. The two
diverge sharply. For a binary indicator the information is $\lambda^2 p(1-p)$, so a
symptom flag endorsed by a few percent of patients contributes almost nothing to
localization however strongly it loads: at $30\%$ endorsement a unit-loading item
carries information $\approx0.10$, at $50\%$ it carries $0.33$, and only a
near-ubiquitous item approaches the information a single well-scaled continuous
measure supplies. The practical consequence is a target-driven item count: to reach
a posterior standard deviation of $0.5$ on an axis requires accumulated information
$\gtrsim3$, which is roughly four items at loading $0.7$ but $\sim\!9$ at loading
$0.5$ and $\sim\!31$ at loading $0.3$. What localizes a patient is the information
an item carries, and information --- not loading --- is the currency the rest of the
instrument spends.

\begin{figure}[t]
  \centering
  \includegraphics[width=\linewidth]{fig3_adaptive.png}
  \caption{\textbf{Adaptive assessment --- the fixed order is a tight upper bound,
  and two axes are bank-limited.} \textbf{(A)} Reliability versus number of items
  administered, most-informative-first, for all eight axes; three to six items
  reach the useful range on well-populated axes, while mania/activation (reliability
  $0.41$, $2$-item bank) and substance ($0.43$, $4$-item bank) are capped by
  instrument content, not missingness. \textbf{(B)} A simulated computerized
  adaptive test (500 synthetic patients per axis, seed $20240704$) versus the
  fixed population-mean order. On Gaussian-topped axes (e.g.\ immunometabolic) the
  two coincide, because $\theta$-independent information leaves nothing to adapt to;
  only Bernoulli-topped axes (e.g.\ suicidality) show a small adaptive signature.
  The population-mean reliability curve is therefore a tight upper bound on adaptive
  efficiency.}
  \label{fig:adaptive}
\end{figure}

\subsection{Adaptive sampling: a tight upper bound, and where the bank is the limit}
\label{sec:res-adaptive}

Accumulating information most-informative-first yields the reliability-versus-items
curve for each axis (Fig.~\ref{fig:adaptive}A). A small number of well-chosen
indicators recovers most of an axis's reliable signal: three to six items reach the
useful range on the well-populated axes. Two axes cannot be rescued by more items
of the same kind --- mania/activation plateaus at reliability $0.41$ (a $2$-item
bank) and substance at $0.43$ (a $4$-item bank). These ceilings are a property of
the \emph{instrument content}, not of missingness: no adaptive schedule and no
completed data matrix can exceed them, because the information simply is not in the
bank. Naming them is a service the instrument performs that a point-estimate map
cannot.

The fixed most-informative-first order is the population-mean curve, an upper bound
on any patient-adaptive schedule. Simulating a full CAT
(\S\ref{sec:m-voi}) shows the bound is \emph{tight} (Fig.~\ref{fig:adaptive}B): on
every axis whose leading items are Gaussian --- whose information is
$\theta$-independent --- patient-adaptive selection coincides with the fixed order
to within $2\times10^{-4}$ in reliability, because there is nothing to adapt to.
Only the two axes with Bernoulli-topped banks show any adaptive signature at all,
and it is small: suicidality diverges by at most $0.012$ (at two items) and
substance by $0.0015$ at its four-item cap. The operational message is
counter-intuitive but useful: for this instrument a fixed, pre-specified battery is
essentially as efficient as a fully adaptive one, so the minimal-indicator rule can
be published as a static list rather than requiring adaptive delivery
infrastructure.

\begin{figure}[t]
  \centering
  \includegraphics[width=\linewidth]{fig4_voi.png}
  \caption{\textbf{Value of information --- designing an efficient shared battery
  and locating under-measured cohorts.} The closed-form entropy reduction from each
  candidate item (Eq.~\ref{eq:m-voi}) drives a greedy cross-axis selection: mean
  reliability across the eight axes rises with battery size, reaching $0.70$ for a
  shared $27$-item battery. Applied per cohort, the same accounting shows where an
  axis is under-measured (e.g.\ schizophrenia carries $20.7$ immunometabolic items
  on average versus $34.9$ in bipolar disorder), concentrating the value of
  additional measurement.}
  \label{fig:voi}
\end{figure}

\subsection{Value of information: designing an efficient shared battery}
\label{sec:res-voi}

The closed-form entropy reduction \eqref{eq:m-voi} lets a greedy rule assemble a
battery that maximises information per item across all eight axes at once.
Figure~\ref{fig:voi} shows the cross-axis reliability rising with battery size; a
shared $27$-item battery reaches a mean reliability of $0.70$ across the eight
axes, a substantial reduction from administering every bank in full. The same
accounting, applied within each cohort, localizes where an axis is under-measured:
schizophrenia patients carry far fewer immunometabolic observations than bipolar
patients ($20.7$ versus $34.9$ items on average) and fewer sleep observations
($3.9$ versus $8.4$), so the value of adding those measurements is concentrated in
the schizophrenia cohort. Value of information thus turns the map into a design
tool --- it prescribes not just how few items suffice in general, but which cohort
should be measured more, and on what.

\begin{figure}[t]
  \centering
  \includegraphics[width=\linewidth]{fig5_calibration.png}
  \caption{\textbf{Trust I --- the posterior coordinate is calibrated.} From
  $10{,}000$ synthetic patients with known coordinates and real bootstrapped
  missingness patterns (seed $20260704$), projected through Fisher-scoring EAP
  (\S\ref{sec:m-calib}). \textbf{(A)} Empirical versus nominal credible level: the
  per-axis mean and the joint eight-dimensional ellipsoid both fall on the identity
  line ($0.949$ at nominal $0.95$). \textbf{(B)} $95\%$ interval coverage
  stratified by the number of observed indicators on the axis, with per-bin
  binomial error bars; coverage is flat from prior-dominated ($0$ items, nominal by
  construction) to richly observed ($11+$).}
  \label{fig:calibration}
\end{figure}

\subsection{Trust I --- the coordinate's uncertainty is calibrated}
\label{sec:res-calib}

An instrument that reports intervals is only useful if the intervals are honest.
Simulating $10{,}000$ patients with known coordinates and real missingness patterns
(\S\ref{sec:m-calib}), the credible intervals are well-calibrated across the board
(Fig.~\ref{fig:calibration}). Empirical coverage tracks nominal at every level ---
$0.501$, $0.801$, $0.900$ and $0.949$ against nominal $0.50$, $0.80$, $0.90$ and
$0.95$ --- both averaged over the eight axes and for the joint eight-dimensional
credible ellipsoid ($0.947$ at the $95\%$ level; Monte-Carlo standard error
$0.0022$). Coverage is flat as a function of how much data the patient supplied
(Fig.~\ref{fig:calibration}B): from prior-dominated axes with zero observed items,
through one or two, up to eleven or more, the $95\%$ interval covers at
$0.947$--$0.952$, every bin within its binomial error of nominal. The zero-item bin
recovers nominal coverage by construction --- the projection returns the prior when
there is no data, and the test confirms the machinery does exactly that. The map's
uncertainty is therefore not decorative: an interval labelled $95\%$ contains the
truth $95\%$ of the time, whether the patient was richly or barely assessed.

\begin{figure}[t]
  \centering
  \includegraphics[width=\linewidth]{fig6_archetype.png}
  \caption{\textbf{Trust II --- the archetype simplex summarizes the continuum, and
  its reconstruction fidelity is quantified.} \textbf{(A)} The five-corner
  archetypal simplex (Eq.~\ref{eq:m-archetype}) with all $9{,}013$ patients
  projected into it, coloured by dominant corner; the interior is densely
  populated. \textbf{(B)} Distribution of archetype-weight entropy: most patients
  are blends, not types. \textbf{(C)} Reconstruction fidelity per axis --- the
  five-corner convex blend ($R^2=0.59$ pooled) against a five-component PCA linear
  optimum ($0.80$) and a five-centroid hard partition ($0.33$); faithful where
  corners separate (immunometabolic $0.89$), weak on the bank-limited axes
  (cognition, substance). \textbf{(D)} Reconstruction error versus archetype
  entropy: ambiguous ``fog'' patients do not reconstruct worse ($r=-0.13$).}
  \label{fig:archetype}
\end{figure}

\begin{figure}[t]
  \centering
  \includegraphics[width=\linewidth]{explain_pca_vs_archetype.png}
  \caption{\textbf{What each low-dimensional summary produces on the map.}
  \textbf{(A)} The five leading principal components: each is a signed
  \emph{contrast} that mixes axes (PC1 loads jointly on developmental risk, sleep,
  severity and suicidality; PC2 is essentially $-$immunometabolic), captures a
  named fraction of variance, and describes a \emph{direction} no patient can
  \emph{be}. \textbf{(B)} The five archetype corners: each is a near-real extreme
  \emph{profile} on the eight axes with a clinical label (C3 $=$ immunometabolic
  $+3.5$, severity $+2.4$; C1 $=$ sleep $+2.8$, mania $+1.7$), so a patient can be
  described as a convex blend --- ``$70\%$ C3'' --- with an attendant clinical
  reading. Both summarize the same $8$-dimensional cloud in five numbers; only the
  archetype corners are nameable poles.}
  \label{fig:pcaarch}
\end{figure}

\subsection{Trust II --- the archetype summary reconstructs the coordinate faithfully}
\label{sec:res-arch}

\emph{The population is a continuum, not a set of types}, and this dictates the
summary. The best five-way partition of the real coordinate cloud separates
patients no better than a structureless Gaussian null (silhouette $0.140$ versus
$0.137\pm0.002$; $z=1.13$, not significant), and most patients sit in the interior
of the simplex as blends rather than at a corner (Fig.~\ref{fig:archetype}B): the
median archetype-weight entropy is well away from zero, and the majority have no
single dominant corner. When there are no natural clusters, a hard-partition
``biotype'' claim is unjustified, and the honest low-dimensional summary is not a
set of groups but a set of \emph{reference poles} that each patient is a graded
mixture of --- exactly what the five-corner archetypal simplex provides
(Fig.~\ref{fig:archetype}A), and exactly what neither a $k$-means partition nor a
principal-component projection can.

\emph{Reconstruction fidelity is the wrong yardstick for this summary, and it is
instructive to see why.} Reconstructing each patient from their five archetype
weights recovers $59.0\%$ of the pooled coordinate variance
(Fig.~\ref{fig:archetype}C). Taken alone that looks modest against a five-component
principal-component projection, the unconstrained linear optimum at the same
dimension, which recovers $79.6\%$. But the map's coordinate variance is nearly
isotropic --- the leading principal component carries only $26.9\%$, and five
components are needed to reach $80\%$ --- so variance capture is a weak test of any
five-dimensional summary here. The expected variance a \emph{random} five-dimensional
subspace retains is exactly $k/d=5/8=0.625$ regardless of the covariance spectrum,
and a $500$-seed simulation confirms it ($0.624\pm0.042$). The archetype simplex
therefore does not merely trail the linear optimum; it sits at the $21$st percentile
of \emph{random} five-dimensional projections --- slightly below what an arbitrary
subspace achieves. Reconstruction $R^2$ cannot be the reason to prefer the simplex,
because a coin-flip subspace beats it, and the fidelity rises only slowly with the
number of corners ($R^2\approx0.53$, $0.62$, $0.69$ at four, five and six archetypes
in a fresh refit whose five-corner value reproduces the production $0.59$ to within
optimizer reinitialization noise; Fig.~\ref{fig:added}B inset). What the simplex provides is what none of these linear
or random summaries do: convex, non-negative, sum-to-one membership in five
\emph{named} clinical poles (Fig.~\ref{fig:pcaarch}) --- a principal component is a
signed contrast that mixes axes and that no patient can \emph{be}, whereas an
archetype corner is a near-real extreme profile a patient can be $70\%$ of, with the
attendant clinical reading. The categorical resolution bears this out: the archetype
assignment accounts for $\eta^2=0.256$ of the variance in the map's coordinates,
against $0.026$ for the DSM-5 diagnosis and $0.018$ for cohort --- an
$\approx\!9.7$-fold tighter description of where a patient sits, on the map's own
axes.

Read this way, the per-axis reconstruction pattern is a diagnostic rather than a
scorecard: the simplex is faithful where the corners are well separated
(immunometabolic $R^2=0.89$, overall severity $0.86$, sleep $0.76$, developmental
risk $0.75$) and near-flat where they are not (cognition $0.06$, substance
$\approx0$) --- the same two bank-limited axes the information analysis already
flagged, now visible as axes along which no five poles could separate patients
because the items carry too little information. And reconstruction does not degrade
for the ambiguous ``fog'' patients: error correlates \emph{negatively} with
archetype entropy (Fig.~\ref{fig:archetype}D; $r=-0.13$, $p<10^{-30}$), because
interior blends lie inside the corner hull whereas near-corner patients include
hull-exceeding extremes. The simplex is thus the defensible summary of a continuum
with no clusters: interpretable poles, graded membership, and a fidelity profile
that reports --- rather than hides --- the axes on which the instrument itself is
weak.

\begin{figure}[t]
  \centering
  \includegraphics[width=\linewidth]{fig7_added_value.png}
  \caption{\textbf{The instrument's choices change decisions.} \textbf{(A)} Building
  the shared battery by ranking items on Fisher information (this paper) versus on
  loading magnitude, over all eight axes: at the $27$-item battery the
  information-ranked rule reaches mean reliability $0.71$ against $0.60$ for the
  loading-ranked rule (a $+0.11$ gap), because high-loading, low-prevalence items
  localize no one. \textbf{(B)} Reconstruction $R^2$ of the five-corner archetype
  simplex placed against the full $0\to1$ range: the raw eight-dimensional
  coordinate ($1.00$), the best five-dimensional linear projection (PCA-5, $0.80$),
  a \emph{random} five-dimensional subspace (mean $0.62$, the analytic $k/d$ value,
  $500$ seeds), the archetype simplex ($0.59$), a five-centroid hard partition
  ($0.33$) and a single centroid ($0.00$). Because the coordinate variance is nearly
  isotropic, a random subspace retains \emph{more} variance than the archetypes;
  $R^2$ is therefore uninformative here, and the case for the simplex rests on
  interpretability, not fidelity. Inset: reconstruction $R^2$ is stable across the
  number of corners.}
  \label{fig:added}
\end{figure}

\subsection{Utility I --- the information rule changes which items you would collect}
\label{sec:res-utility-battery}

Two of the instrument's claims --- that information, not loading, should rank items,
and that the archetype simplex is worth its fidelity cost --- can be tested as
head-to-head design decisions rather than asserted. Both hold, and both matter
(Fig.~\ref{fig:added}). Assembling the shared battery by the information rule of
\S\ref{sec:res-voi} against the naive alternative of ranking items by loading
magnitude, at the same $27$-item budget, yields mean reliability $0.705$ for the
information rule against $0.599$ for the loading rule --- a $0.106$ gap across the
eight axes (Fig.~\ref{fig:added}A). The loading rule spends its early picks on
high-loading but low-prevalence indicators that localize almost no one, and does not
catch up until roughly item sixteen. Choosing items by the information they carry is
thus not a refinement of choosing by loading; it selects a materially different, and
measurably better, battery.

The archetype summary invites the opposite test --- whether its interpretability is
bought at an indefensible fidelity cost --- and the answer reframes what the
$59\%$ reconstruction figure means (Fig.~\ref{fig:added}B). Anchored against the
full range, the archetype simplex ($0.59$) sits below not only the five-component
linear optimum ($0.80$) but a \emph{random} five-dimensional subspace ($0.62$),
because the near-isotropic coordinate variance makes any five-dimensional projection
capture roughly $k/d=0.625$. Reconstruction $R^2$ cannot adjudicate between the
summaries, and the choice is therefore made on what only the simplex delivers ---
named, convex, per-patient membership in five clinical poles --- not on variance
retained.

\begin{figure}[t]
  \centering
  \includegraphics[width=\linewidth]{fig8_blockers.png}
  \caption{\textbf{The posterior earns its keep, and its intervals degrade
  gracefully.} \textbf{(A)} Head-to-head prognosis on model-simulated data with a
  known remission rule: the full expected-a-posteriori coordinate (this paper)
  against a naive per-axis sum-score, as item completeness falls. The posterior wins
  at every level and the gap \emph{widens} as data get sparse ($+0.02$ at full data,
  $+0.04$ at $25\%$ observed), because it shrinks noisy inputs toward the prior
  instead of averaging them blindly; over $12$ seeds the posterior wins $100\%$ of the
  time. The dotted line is the oracle AUC on the true latent ($0.80$).
  \textbf{(B)} Coverage of the nominal $95\%$ interval when the data-generating
  process \emph{violates} a model assumption, scored with the unperturbed
  projection: correlated residuals (local-independence violation), an omitted ninth
  factor, mismatched loadings, and heavy-tailed residuals. Coverage degrades
  smoothly and never collapses --- the worst case, an extreme local-dependence
  violation, still covers at $0.89$, and heavy tails are tolerated essentially
  without loss.}
  \label{fig:blockers}
\end{figure}

\subsection{Utility II --- carrying the posterior improves downstream prognosis, and calibration is robust}
\label{sec:res-utility-prognosis}

The instrument's foundational promise is that a coordinate carried with its
uncertainty is worth more than a point estimate. Simulating item responses from the
fitted model with a known latent-to-remission rule, and comparing the full
expected-a-posteriori coordinate against a naive per-axis sum-score of the same
observed items, the posterior predicts the outcome better at every level of item
completeness, and the advantage grows as data thin
(Fig.~\ref{fig:blockers}A): the out-of-sample AUC gap is $+0.018$ on complete data
and $+0.039$ at a quarter of items observed, against an oracle AUC of $0.799$ on the
true latent. The mechanism is the posterior's shrinkage --- where a sum-score
propagates the full noise of a sparsely-measured axis, the projection pulls it
toward the prior in proportion to how little was observed --- which is precisely why
the gain concentrates in the sparse regime. The effect is not a simulation artifact:
over twelve seeds and an alternative outcome rule the posterior wins every time (gap
$+0.019\pm0.003$ on full data, $+0.032\pm0.006$ when sparse). We note the boundary of
this claim honestly: on the real cohort, triaging patients by their posterior
uncertainty does not raise remission prediction (out-of-sample AUC $0.709$ whether
or not the least-certain patients are set aside), because remission at follow-up is
driven substantially by treatment and course factors the baseline coordinate cannot
see; the demonstrated value of the posterior is the sharper coordinate, not a
confidence-triage of who is predictable.

Because the calibration of \S\ref{sec:res-calib} samples from the model's own
parameters, it establishes internal honesty but not robustness to misspecification.
We therefore generate data from four data-generating processes that each violate a
model assumption --- correlated residuals (breaking local independence), an omitted
ninth latent factor, multiplicatively mismatched loadings, and heavy-tailed
(Student-$t$) residuals --- and score them with the \emph{unperturbed} projection
(Fig.~\ref{fig:blockers}B). Coverage of the nominal $95\%$ interval degrades
gracefully in every case and collapses in none. Heavy tails are absorbed almost
without loss (coverage $0.955$ at $t$ with three degrees of freedom); an omitted
factor and mismatched loadings sag gently to $0.93$--$0.94$; the steepest and most
consequential violation, correlated residuals, drifts monotonically to $0.889$ at an
extreme shared-nuisance loading --- the expected direction, since local dependence is
where factor-model intervals are most optimistic, yet even there the interval loses
only six points of coverage. Calibration under this instrument is thus not a
by-construction artifact: the intervals are honest when the model holds and fail
predictably, not catastrophically, when it does not, with local dependence flagged
as the assumption most worth checking on real data.
